\subsection{Puertas de Pauli}
	Las conocidas puertas de Pauli $\gate{I}, \gate{X}, \gate{Y}, \gate{Z}$ pueden expresarse de acuerdo con la notación introducida anteriormente como sigue:
	\begin{align}
		\gate{I}_k\ket{a} &= a_{k0}\stateket{a}{k}{0}+a_{k1}\stateket{a}{k}{1} \\
		\gate{X}_k\ket{a} &= a_{k0}\stateket{a}{k}{1}+a_{k1}\stateket{a}{k}{0}\\
		\gate{Y}_k\ket{a} &= ia_{k0}\stateket{a}{k}{1}-i a_{k1}\stateket{a}{k}{0}\\
		\gate{Z}_k\ket{a} &= a_{k0}\stateket{a}{k}{0}-a_{k1}\stateket{a}{k}{1}
	\end{align}

\subsection{Puertas de rotación de Pauli}
Las tres puertas de rotación de Pauli $\gate{RX}, \gate{RY}$ y $\gate{RZ}$ se definen tomando exponenciales de la puerta de Pauli correspondiente.

Para un valor $\theta\in\R$:
\begin{align}
    \gate{RX}(\theta)_k\ket{a} &= \left(a_{k0}\cos(\frac{\theta}{2})-a_{k1}\sin(\frac{\theta}{2})i\right)\stateket{a}{k}{0}\nonumber\\
    &\quad+\left(a_{k1}\cos(\frac{\theta}{2})-a_{k0}\sin(\frac{\theta}{2})i\right)\stateket{a}{k}{1}\\
    \gate{RY}(\theta)_k\ket{a} &= \left(a_{k0}\cos(\frac{\theta}{2})-a_{k1}\sin(\frac{\theta}{2})\right)\stateket{a}{k}{0}\nonumber\\
    &\quad+\left(a_{k1}\cos(\frac{\theta}{2})+a_{k0}\sin(\frac{\theta}{2})\right)\stateket{a}{k}{1}\\
    \gate{RZ}(\theta)_k\ket{a} &= a_{k0}\left(\cos(\frac{\theta}{2})-\sin(\frac{\theta}{2})i\right)\stateket{a}{k}{0}\nonumber\\
    &\quad+a_{k1}\left(\cos(\frac{\theta}{2})+\sin(\frac{\theta}{2})i\right)\stateket{a}{k}{1}
\end{align}

\begin{proposition}
	Sea $A$ una de las tres matrices de Pauli $X$, $Y$ o $Z$. Sea $k\in\N_n$ y $\theta, \varphi\in\R$.
	\begin{equation}
		\gate{RA}(\theta)_k\gate{RA}(\varphi)_k=\gate{RA}(\theta+\varphi)_k
	\end{equation}
\end{proposition}
